% !TeX encoding = UTF-8
% !TeX spellcheck = pt_BR
% !TeX root = 20201-CODDISC-JOAO-FERREIRA.tex
% ------------------------------------------------------------------------------

% ------------------------------------------------------------------------------
% Template para aulas da UNINASSAU
% Autor: João Ferreira da Silva Júnior
% Versão: 0.2 - 2021-09-12
% ------------------------------------------------------------------------------



% ------------------------------------------------------------------------------
% ------------------------------------------------------------------------------
% ------------------------------------------------------------------------------
% ------------------------------------------------------------------------------
% ------------------------------------------------------------------------------
% ------------------------------------------------------------------------------
% ------------------------------------------------------------------------------
% ------------------------------------------------------------------------------
% ------------------------------------------------------------------------------
% ------------------------------------------------------------------------------
\section[Plano de Ensino]{Plano de Ensino}\label{sec:sec02}


{\usebackgroundtemplate{{\transparent{0.2}\includegraphics[width=\paperwidth,height=\paperheight,keepaspectratio]{img/logo/background-branco.png}}}
% \begin{frame}[allowframebreaks,t]%\frametitle{ ---}
  \begin{frame}[plain]%\frametitle{}

    \vfill
    \centering

    \addtobeamertemplate{block begin}{\pgfsetfillopacity{0.5}}{\pgfsetfillopacity{1}}
    % \setbeamercolor{block title}{use=structure,fg=white,bg=blue}
    \setbeamercolor{block body}{use=structure,fg=black,bg=blue!50!white}
    \begin{block}{}
      \centering{}
      \Huge{\textbf{Plano de Ensino}}
    \end{block}

    \vfill

\end{frame}
} % usebackgroundtemplate



% ------------------------------------------------------------------------------
% ------------------------------------------------------------------------------
% ------------------------------------------------------------------------------
% ------------------------------------------------------------------------------
% ------------------------------------------------------------------------------
\subsection[Ementa]{Ementa}\label{subsec:sec02-ementa}



% \begin{frame}[allowframebreaks,t]\frametitle{Ementa ---}
\begin{frame}[t]\frametitle{Ementa}

  \begin{itemize}
    \justifying{}
    \setlength\itemsep{1em}
    \item Item 1.
    \item Item 2.
    \item Item 3.
  \end{itemize}

\end{frame}



% ------------------------------------------------------------------------------
% ------------------------------------------------------------------------------
% ------------------------------------------------------------------------------
% ------------------------------------------------------------------------------
% ------------------------------------------------------------------------------
\subsection[Conteúdo Programático]{Conteúdo Programático}\label{subsec:sec02-conteudo-programatico}



% \begin{frame}[allowframebreaks,t]\frametitle{Conteúdo Programático ---}
\begin{frame}[t]\frametitle{Conteúdo Programático}

  \begin{itemize}
    \justifying{}
    \setlength\itemsep{1em}
    \item Item 1.
    \item Item 2.
    \item Item 3.
  \end{itemize}

\end{frame}



% ------------------------------------------------------------------------------
% ------------------------------------------------------------------------------
% ------------------------------------------------------------------------------
% ------------------------------------------------------------------------------
% ------------------------------------------------------------------------------
\subsection[Metodologia de Ensino]{Metodologia de Ensino}\label{subsec:sec02-metodologia-ensino}



% \begin{frame}[allowframebreaks,t]\frametitle{Metodologia de Ensino ---}
\begin{frame}[t]\frametitle{Metodologia de Ensino}

  \begin{block}{}
    \justifying{}
    A disciplina, dependendo de sua natureza, pode ser ministrada através de conteúdos teóricos, conteúdos práticos, aulas de campo em instituições específicas e ainda pode utilizar recursos de exposições dialogadas, grupos de discussão, seminários, debates competitivos, apresentação e discussão de filmes e casos práticos, onde os conteúdos podem ser trabalhados mais dinamicamente, estimulando o senso crítico e científico dos alunos.
  \end{block}

\end{frame}



% ------------------------------------------------------------------------------
% ------------------------------------------------------------------------------
% ------------------------------------------------------------------------------
% ------------------------------------------------------------------------------
% ------------------------------------------------------------------------------
\subsection[Metodologia de Avaliação]{Metodologia de Avaliação}\label{subsec:sec02-metodologia-avaliacao}



% \begin{frame}[allowframebreaks,t]\frametitle{Metodologia de Avaliação ---}
\begin{frame}[t]\frametitle{Metodologia de Avaliação}

  \begin{itemize}
    \justifying{}
    \setlength\itemsep{1em}
    \item No decorrer de cada período letivo são desenvolvidas 02 (duas) avaliações por disciplina, para efeito do cálculo da média parcial.
    \item A média parcial é calculada pela média aritmética das duas avaliações efetuadas. O aluno que alcançar a média parcial maior ou igual a 7,0 (sete) é considerado aprovado.
    \item O aluno que não alcançar a média parcial faz em exame final onde precisa alcançar média final maior ou igual a 5,0.
    \item São aplicadas avaliações dos tipos: provas teóricas, provas práticas, seminários, trabalhos individuais ou em grupo e outras atividades em classe e extraclasse.
    \item O exame final é, obrigatoriamente, prova escrita.
  \end{itemize}

\end{frame}



% ------------------------------------------------------------------------------
% ------------------------------------------------------------------------------
% ------------------------------------------------------------------------------
% ------------------------------------------------------------------------------
% ------------------------------------------------------------------------------
\subsection[Bibliografia Sugerida]{Bibliografia Sugerida}\label{subsec:sec02-bibliografia-sugerida}



% \begin{frame}[allowframebreaks,t]\frametitle{Bibliografia Básica ---}
\begin{frame}[t]\frametitle{Bibliografia Básica}

  Exemplos de citações:
  \begin{itemize}
    \justifying{}
    \setlength\itemsep{1em}
    \item TANENBAUM, A. S. Redes de Computadores. 5o ed. Pearson Prentice Hall, 2011. \cite{Tanenbaum2011}
    \item TANENBAUM, A. S.; BOS, H. Sistemas Operacionais Modernos. 4o ed. Pearson Prentice Hall, 2016. \cite{Tanenbaum2016}
  \end{itemize}

\end{frame}



% \begin{frame}[allowframebreaks,t]\frametitle{Bibliografia Complementar ---}
\begin{frame}[t]\frametitle{Bibliografia Complementar}

  Exemplos de citações:
  \begin{itemize}
    \justifying{}
    \setlength\itemsep{1em}
    \item TANENBAUM, A. S. Redes de Computadores. 5o ed. Pearson Prentice Hall, 2011. \cite{Tanenbaum2011}
    \item TANENBAUM, A. S.; BOS, H. Sistemas Operacionais Modernos. 4o ed. Pearson Prentice Hall, 2016. \cite{Tanenbaum2016}
  \end{itemize}

\end{frame}
